
Within this reduced spatial requirement, channel selection was stable at the level of scalp region but not at the level of a single electrode. Fp1 and Fp2 accounted for the large majority of selections in both corpora: Raja favoured Fp2 in 75/184 selections and Fp1 in 67/184, whereas Cao2018 favoured Fp1 in 119/232 and Fp2 in 57/232. However, the four conditions selected the same channel in only 9/46 Raja sessions and 6/58 Cao2018 sessions, with pooled full agreement of 15/104. This apparent tension indicates that the detector consistently localised blink information to the frontopolar region while the electrode carrying the strongest signal varied across conditions and sessions. Such variation is consistent with lateralised ocular projections and inter-subject differences in blink topography, electrode placement, and signal amplitude rather than with instability of the underlying detection principle~\citep{plochl2012combining,guttmann2019new}. The practical implication is that deployment should retain both frontopolar sites and select between them adaptively, rather than fixing acquisition or detection to Fp1 or Fp2 alone.
